% Section on the 2022 music census microdata. This is the only current
% Austin-specific measurement of this workforce, and its retention finding is
% the report's most novel result.
\section{What musicians say about staying}
\label{sec:census}

Section~\ref{sec:venues} counted rooms. The only current Austin-specific
measurement of the people who play in those rooms is the 2022 Greater Austin
Music Census, an online instrument that Sound Music Cities fielded with the
City of Austin Music \& Entertainment Division between 15 July and 12 September
2022.\footnote{Sound Music Cities is the consultancy that has run the same
instrument across a rolling cohort of U.S.\ cities since 2022 (Sacramento,
Nashville, New Orleans, Washington D.C., Minneapolis, Cleveland, Baltimore,
and others; the full roster is listed in the references). The Austin fielding
was distributed through community mailing lists, the city's Music \&
Entertainment Division channels, and partner organizations. The city
publishes the individual respondent records as an open dataset,
\textit{2022 Austin Music Census: Details}, Socrata identifier
\texttt{an3p-3yqx}, at \url{https://data.austintexas.gov/resource/an3p-3yqx.csv},
retrieved 1~August 2026. The three Sound Music Cities reports that accompany
the microdata (\textit{Summary Report}, \textit{Data Appendix}, and
\textit{DEI Data Appendix}) are the source for anything printed above the
respondent level; the underlying \NCensusNRespondents\ records are the source
for this study. The Summary Report headline count is 2,260; the microdata
release carries 33 fewer rows and the release documents no reason. Every share
in this section uses the \NCensusNRespondents-row microdata file, so a value
here can sit a few tenths of a point away from a published percentage even
where the two agree.} The 2022 instrument dropped the dollar-denominated income
questions the 2014 fielding had asked, so it cannot say what a musician earns,
but it kept a great deal else: work volume, workspace, tenure, insurance,
business spending, three-year intent, and a ranked list of business
pressures.\footnote{The 2014 census, run by the same authoring team on behalf
of the city, asked for calendar-year 2013 music income in banded form. The
2022 authors state that they dropped the item to raise completion rates and
because respondents had told them the frame did not fit self-employed and
mixed-income careers; see \textit{Summary Report}, printed p.~6, at
\url{https://www.austinmusiccensus.org/}. The consequence is that Austin's own
music census has not asked what a musician earns for the twelve years running.}
This study reads the census from the respondent-level file rather than the
published tables, because a per-respondent construction is needed for two
questions the printed report cannot answer: whether one person's two three-year
answers move together, and whether the widely cited spending headline describes
any actual person.

That file records \NCensusNRespondents\ respondents across 107 variables. Every
number in this section describes those respondents and not the Austin music
workforce.\footnote{The census was a self-selected online instrument with no
design or post-stratification weights. It is not a probability sample of
Austin's music workforce, and a probability-sample benchmark for the same
population does not exist. Every share in this section is therefore a
descriptive statistic on the people who answered, not a population estimate,
and none is comparable with the American Community Survey estimates in
Section~\ref{sec:earnings}, the Bureau of Labor Statistics wage series in
Section~\ref{sec:payroll}, or the Comptroller receipts panel in
Section~\ref{sec:venues}. No standard error or margin of error is computed
anywhere in this section. Item nonresponse is heavy and uneven, and several
blocks were routed only to performers or only to venue operators, so every
statistic in this section is reported against its own denominator. Six
figures computed here reproduce the Sound Music Cities \textit{Data Appendix}
to within a percentage point (the two three-year retention shares, the mean
and median distance travelled to work, the sum-of-means spending headline, and
the local share of spending), which is the best available check that the file
is being read correctly.}

\subsection{Most respondents expect to stay in music; far fewer expect to stay in Austin}

The same \NCensusRetentionDenom\ respondents answered both three-year questions,
so the comparison below is between one person's two answers rather than between
two groups. \NCensusIntendContinueMusic\ answered definitely or maybe yes to
still working in music, and \NCensusIntendStayAustin\ answered definitely or
maybe yes to still living in greater Austin, a within-person gap of
\NCensusRetentionGap.\footnote{Each share is an unweighted count over its own
denominator, \(\hat p = 100 \times n_{\mathrm{yes}}/n_{\mathrm{answered}}\),
where \(n_{\mathrm{yes}}\) counts respondents answering definitely or maybe yes
and \(n_{\mathrm{answered}}\) counts everyone who answered that item. A
respondent who skipped the item is out of both numerator and denominator
rather than counted as a no. The retention gap is the difference of the two
shares on the same respondents,
\(\hat p_{\mathrm{music}} - \hat p_{\mathrm{Austin}}\), which is a within-person
comparison and not a comparison of two groups. No weight enters anywhere in
this section, because the release carries none, so no margin of error can be
computed. Both retention shares reproduce the Sound Music Cities
\textit{Summary Report} p.~7 to the tenth of a percentage point on the
microdata file.} Most of the complement to the Austin share is uncertainty
rather than stated departure: among the \NCensusRetentionDenom\ respondents,
about a quarter answered ``unsure'' and roughly one in ten answered maybe or
definitely no, which is a weaker claim than a third planning to
leave.\footnote{The Austin item's five options are definitely yes, maybe yes,
unsure, maybe no, and definitely no. The full breakdown on the
\NCensusRetentionDenom\ respondents is 44.2\% definitely yes, 19.9\% maybe
yes, 25.6\% unsure, 6.7\% maybe no, and 3.6\% definitely no; the two shares
combined into ``definitely or maybe yes'' produce \NCensusIntendStayAustin.
Read the 25.6\% unsure share as respondents who could not answer, not as
respondents who leaned toward leaving.}

\begin{figure}[htbp]
\centering
\includegraphics[width=\textwidth]{fig26_retention_split.png}
\caption{Share expecting to continue music work and to remain in greater Austin
three years on, overall and by housing tenure. Source: 2022 Greater Austin
Music Census respondent-level file, Sound Music Cities and City of Austin
Music \& Entertainment Division, City of Austin open data dataset
\texttt{an3p-3yqx}, retrieved 1~August 2026. How to read it: both questions
were answered by the same \NCensusRetentionDenom\ respondents, so the distance
between the two bars is a within-person gap and not a comparison of two
groups. Caveats: this is a self-selected online survey with no weights, so
the percentages describe respondents rather than the Austin music workforce,
and they are not comparable with the American Community Survey estimates in
Section~\ref{sec:earnings}.}
\label{fig:retention}
\end{figure}

\subsection{Housing tenure separates the answers most sharply; sector and workspace do not separate them at all}

Who are the doubtful respondents? Of six respondent characteristics this study
tested against the Austin-retention answer, housing tenure separates it most
sharply: \NCensusStayHousingTenureHomeowner\ of homeowners expected to stay in
greater Austin against \NCensusStayHousingTenureRenter\ of renters, with
\NCensusStayHousingTenureOtherOrUnstableHousing\ of respondents in other or
unstable housing situated between them.\footnote{The homeowner--renter
difference in expected-stay rates is unlikely to be chance among these
respondents: \NCensusStaychiTwoHousingTenure. The test is Pearson's
chi-squared for independence in a two-way table,
\(X^{2} = \sum_{j}\sum_{k}(O_{jk} - E_{jk})^{2}/E_{jk}\) with
\(E_{jk} = R_j C_k / N\), where \(O_{jk}\) is the number of respondents in row
\(j\) and column \(k\), \(R_j\) and \(C_k\) are the row and column totals,
\(N\) is the number answering both items, and \(E_{jk}\) is the count expected
if the two answers were unrelated. The number in parentheses after the
statistic is the degrees of freedom, \((J-1)(K-1)\) for a table with \(J\)
rows and \(K\) columns. In a logistic regression conditioning on primary
music-ecosystem sector, years of experience, and health insurance, renters
have \NCensusLogitRenterOr\ times the odds of expecting to stay in Austin
that homeowners do, on \NCensusLogitN\ respondents who answered every item in
the equation. The specification is
\[
\ln\frac{\Pr(\mathrm{stay}_i = 1)}{1 - \Pr(\mathrm{stay}_i = 1)}
  = \alpha + \sum_s \beta_s\,\mathrm{sector}_{si}
  + \sum_e \gamma_e\,\mathrm{exper}_{ei}
  + \sum_h \delta_h\,\mathrm{tenure}_{hi}
  + \theta\,\mathrm{insured}_i ,
\]
where \(\mathrm{stay}_i\) is one when respondent \(i\) answered definitely or
maybe yes to still living in greater Austin in three years, the
\(\mathrm{sector}_{si}\) are indicators for primary music-ecosystem sector,
the \(\mathrm{exper}_{ei}\) for years-of-experience band, the
\(\mathrm{tenure}_{hi}\) for housing tenure with homeowner omitted, and
\(\mathrm{insured}_i\) is one for holding health insurance. The quantity
quoted is the odds ratio \(e^{\hat\delta_{\mathrm{renter}}}\); a value below
one means renters had lower odds of expecting to stay than homeowners who
match them on sector, experience, and insurance. An odds ratio is not a ratio
of probabilities and sits further from one than the probability ratio does
when the outcome is common, as it is here. Estimation is by maximum
likelihood on the respondents who answered every item in the equation, which
is fewer than answered the retention question alone. Neither the test nor the
regression is weighted, because the census carries no design weights, and
both p-values assume a simple random sample, which a self-selected online
survey is not. Read the test statistics as descriptive discriminants among
respondents rather than as inference to the Austin music workforce.
Tenure was not assigned at random, and income, age, and family stage all
move with it, none of which the census measured; the tenure result is
therefore an association among respondents and not the cost of renting.
The analysis recodes tenure into three mutually exclusive categories, because
the published appendix's categories overlap, which is why the tenure shares
here do not match the \textit{Data Appendix}.} Health insurance separates the
answers the same way: \NCensusStayHealthInsuranceHasHealthInsurance\ of the
1,390 respondents with health insurance expected to stay against
\NCensusStayHealthInsuranceNoHealthInsurance\ of the 254
uninsured.\footnote{The census item is respondent-reported coverage of any
kind (an employer plan, a marketplace plan, Medicaid, Medicare, VA, or
another source), asked as a single yes/no. It is not directly comparable with
the American Community Survey health-insurance coverage rate quoted in the
summary, which uses the ACS multi-item HICOV construction on a probability
sample of households.}

Two music-career characteristics also separate the Austin-retention answer,
though not in a straight line. Years of experience is the second strongest of
the six discriminants (\NCensusStaychiTwoYearsOfExperience), and it runs
non-monotonically across the four experience bands the census reports: the
lowest expected-stay rate belongs to the 6-to-10-year band at
\NCensusStayYearsOfExperienceSixToOneZeroYears, below both newer entrants at
\NCensusStayYearsOfExperienceUnderThreeYears\ (under three years of
experience) and the most experienced at
\NCensusStayYearsOfExperienceMoreThanOneZeroYears\ (more than ten years). How
much of a respondent's music income comes from local gigs also separates the
Austin answer, on the same 1,900 respondents
(\NCensusStaychiTwoMusicIncomeFromLocalGigs).

Two other characteristics do not separate the answers. Neither the
respondent's primary music-ecosystem sector (music creative, music industry,
or venue or presenter) nor the respondent's workspace status (owns, rents,
provided or not needed, or needs but lacks) discriminates the Austin-retention
answer at conventional levels; both are within the range chance would
produce among respondents
(\NCensusStaychiTwoPrimarySector; \NCensusStaychiTwoWorkSpaceStatus).

\begin{keyfinding}
{\small That result ties together two parts of this study.
Section~\ref{sec:costs} showed that Austin's cost problem is centered on
housing more than on prices generally, and this section finds that housing tenure separates the
census respondents' Austin-retention answer more sharply than any of the six
respondent characteristics tested. Health insurance splits the same answer
the same way, so the two strongest household-security measures both
separate stayers from the unsure, while primary sector and workspace access
do not separate them at all.}
\end{keyfinding}

\subsection{Working in music costs money, and the widely cited figure overstates the typical case}

\begin{figure}[htbp]
\centering
\includegraphics[width=\textwidth]{fig27_cost_of_working.png}
\caption{Annual respondent spending on their own music work, by category,
music creatives only, 2025 dollars. Source: 2022 Greater Austin Music Census
respondent-level file (dataset \texttt{an3p-3yqx}), converted from 2022 to
2025 dollars with the national consumer price index (CPI-U annual factor).
How to read it: mean and median are shown for each category because the
distribution is heavily right-skewed and the median is the better guide to a
typical respondent. Caveats: each category has its own base, ranging from 611
to 742 music creatives; the census is a self-selected online sample and
carries no weights.}
\label{fig:costofworking}
\end{figure}

Coverage of the 2022 census in local media and in the Sound Music Cities
\textit{Data Appendix} often cites an annual music-work spending figure of
about \$10,500 in 2022 dollars, worth \NCensusSpendSumofmeansTwoZeroTwoFive\
in 2025 prices.\footnote{The published \$10,500 headline appears in the Sound
Music Cities \textit{Data Appendix}, printed p.~46. It was widely picked up
locally, including in Austin City Council briefings on the Live Music Fund
and in coverage by KUT News and \textit{Austin Monitor}. The 2022 dollar
inflation factor to 2025 is the ratio of the national CPI-U all-urban annual
averages from the Bureau of Labor Statistics, series identifier
\texttt{CUUR0000SA0}, at \url{https://www.bls.gov/cpi/data.htm}; the same
factor is used throughout this study.} That headline adds together eight
category means computed on eight different sets of respondents, so it does
not describe any actual person. Among the \NCensusSpendCompleteN\ music
creatives who answered every category, the median total is
\NCensusSpendTotalMedian\ and the mean is \NCensusSpendTotalMean, both in
2025 dollars, so the correctly constructed mean sits well below the widely
cited figure and the median sits far below it.\footnote{The widely cited
headline is \(\sum_{k=1}^{8}\bar{x}_k\), where \(\bar{x}_k\) is the mean of
spending category \(k\) taken over \(S_k\), the set of respondents who
answered that category. The eight sets are not the same people. The
per-respondent figures quoted here are built the other way round, summing
within a person first and taking the middle of those totals:
\(\tilde{x} = \operatorname{median}_{i \in S}
   \bigl(\sum_{k=1}^{8} x_{ik}\bigr)\)
over \(S = \bigcap_k S_k\), the respondents who answered every category, so
every dollar in the sum belongs to one person. The two constructions agree
only when each category has the same respondents and the within-person total
is symmetrically distributed, and neither condition holds here. Both are
converted from 2022 to 2025 dollars by \(C_{2025}/C_{2022}\), the ratio of
national CPI-U annual averages, as is the published headline quoted above.
The eight categories in the instrument are gear and rentals, new recordings,
publicity and promotion, merchandise, rehearsal or work space, supplies,
accounting and legal, and travel; the item is asked of music creatives only
and does not appear on the venue-operator branch.} For a single number that
describes a typical respondent, prefer the median: \NCensusSpendTotalMedian\
on the \NCensusSpendCompleteN\ complete-case music creatives, or
\NCensusSpendTotalMedianAnyitem\ on the 873 who answered at least one
category (with unanswered categories treated as zero, a deliberate lower
bound because the release cannot distinguish a true zero from a blank).

These self-reported spending figures cannot be netted against the
Nonemployer Statistics receipts figure quoted in
Section~\ref{sec:payroll}.\footnote{U.S.~Census Bureau, Nonemployer
Statistics, at
\url{https://www.census.gov/programs-surveys/nonemployer-statistics.html}.
Nonemployer Statistics is an administrative-record product built from
business income-tax returns of firms with no paid employees and receipts of
at least \$1{,}000; it reports the receipts side, before expenses, and
carries no expense series. The two numbers therefore cannot be differenced.
They cover different populations (music-industry nonemployer firms in the
Nonemployer file, music creatives in the census), different years, and
different units (businesses in one case and people in the other), and the
census figure is a self-reported cost rather than a tax-return revenue
line.}

\subsection{Respondents who present live music name property tax as their leading pressure}

Among the \NCensusNPresentersRanking\ respondents who both present live music
and completed the nine-item business-pressure ranking,
\NCensusPressureTopThreePropertyTax\ placed property tax in their top three
and \NCensusPressureFirstPropertyTax\ ranked it their single greatest
pressure, ahead of talent costs (\NCensusPressureTopThreeTalentCosts\ in the
top three) and labor cost and retention
(\NCensusPressureTopThreeLaborCostAndRetention).\footnote{The ranking item
asked respondents who present live music, whether as a primary or a
secondary activity, to rank nine business pressures from most to least
pressing. The nine items were building operations and upkeep, changing
audience behavior, labor cost and retention, marketing, neighborhood
redevelopment, ordinances and permitting, property tax, talent costs, and
unpredictable operating costs. \NCensusNPresentersRanking\ respondents
completed the full ranking; that is 8~percent of the \NCensusNRespondents\
respondents in the file and the thinnest block used in this section. This is
a presenter-activity population rather than a venue census: only 77 of the
177 named venue or presenter as their primary sector, and the remaining 100
are music creatives and industry professionals who also present gigs. The
Sound Music Cities \textit{Data Appendix} slides~56--57 split operators from
promoters; the promoter cell falls below this study's 50-record publication
floor and is not reported here.}

That ranking and the closure record in Section~\ref{sec:venues}, where
\NCensusClosuresPropertyShare\ of the 36 dated Austin venue closures cite a
real-estate cause (sale, rent increase, lease expiration, or redevelopment)
against 13 that cite the pandemic, both point at property costs. They do not
point at the same property cost: the census item is a tax that falls on
owners, and the Section~\ref{sec:venues} closure record is dominated by
rent, sale, and lease-tenure disputes at what are mostly tenant venues. The
two sources also diverge in one place. Talent costs rank second among
presenter pressures in the census, but appear as the cause of no closure in
the Section~\ref{sec:venues} record. The Comptroller receipts panel used in
Section~\ref{sec:venues} has no field for what a venue pays performers, so
the talent-cost item is untested here rather than contradicted.

Sections~\ref{sec:city} and \ref{sec:state} turn from what these respondents
report to what the city and the state actually spend on this workforce.
